% Generated draft: requires mathematical and source review.
\documentclass[11pt]{article}
\usepackage[margin=25mm]{geometry}
\usepackage{fontspec}
\setmainfont{Latin Modern Roman}
\usepackage{amsmath,amssymb,mathtools,unicode-math}
\setmathfont{Latin Modern Math}
\usepackage{enumitem,needspace,xurl,hyperref}
\setlength{\parindent}{0pt}
\setlength{\parskip}{5pt}
\setlength{\emergencystretch}{3em}
\newcommand{\N}{\mathbb{N}}
\newcommand{\Z}{\mathbb{Z}}
\newcommand{\Q}{\mathbb{Q}}
\newcommand{\R}{\mathbb{R}}
\newcommand{\C}{\mathbb{C}}
\newcommand{\F}{\mathbb{F}}
\newcommand{\A}{\mathbb{A}}
\newcommand{\PP}{\mathbb{P}}
\newcommand{\E}{\mathbb{E}}
\newcommand{\eps}{\varepsilon}
\newcommand{\dd}{\,\mathrm{d}}
\begin{document}
\section*{UnsolvedMath ID: 11100059}
\textbf{Source problem code:} catalogue entry
\subsection{Definitions and mathematical statement}
Exhibit a model category \((\mathcal M,W,\mathrm{Cof},\mathrm{Fib})\) and prove that its given model structure is not cofibrantly generated. In the standard definition, cofibrant generation requires sets \(I,J\) with
\[
 I\text{-inj}=W\cap\mathrm{Fib},\qquad
 J\text{-inj}=\mathrm{Fib},
\]
where inj denotes the right lifting property, together with the smallness conditions needed for the small-object factorizations. The required negative proof must rule out all such generating sets, not merely one proposed choice.

The source proposes a chain-complex model whose weak equivalences are chain homotopy equivalences. In a standard unbounded \(R\)-module version, cofibrations are degreewise split monomorphisms and fibrations degreewise split epimorphisms; a weak equivalence has an inverse up to chain homotopy. The ring, boundedness and exact variant are not given in the source and must be fixed before proving non-generation. Ring dependence is essential: over a field, splitting makes the homotopy-equivalence and quasi-isomorphism notions coincide, so an unqualified claim that every such chain model is non-generated cannot be used.

The source credits Mike Cole with an early written description and speculates Quillen may have known the structure. The summary reports a later non-generation result but offers ambiguous ``and/or'' references and no ring hypotheses; those reports are not themselves a proof for the unspecified candidate.
\subsection{Short English statement}
Give a model category with a proof that no sets of generating cofibrations and trivial cofibrations can define its model structure; the proposed chain-homotopy model needs a specified base ring.
\subsection{Sources}
\begin{itemize}
\item Mark Hovey, \emph{Algebraic Topology Problem List}, ``Model categories,'' item 13, surviving mirror \url{https://www-users.cse.umn.edu/~tlawson/hovey/}. Recorded access 2026-07-29; extraction html-ordered-list. Mike Cole and Quillen are mentioned without full references.
\item Supplied summary: Christensen--Hovey, ``Quillen model structures for relative homological algebra'' (2001); purported Christensen--Raptis ``Six model structures for DG-modules'' and ``Realizing spectra in ...''; and Raptis's thesis. The summary's author/title assignments and alternatives are incomplete and unverified, not reconstructed into a reliable citation.
\item Local UnsolvedMath registry, \texttt{data/unsolvedmath/problems.json}, permanent ID \(11100059\), exact source code \texttt{catalogue entry}; registry open/OPEN-TRIAGE with a solved report in the summary.
\end{itemize}
\end{document}
